\documentclass[12pt]{article}
\usepackage[T1]{fontenc}
\usepackage[utf8]{inputenc}
\usepackage{lmodern}
\usepackage{textcomp}
\usepackage{enumitem}
\usepackage{geometry}
\usepackage{graphicx}
\usepackage{amsmath, amssymb}
\geometry{margin=1in}
\begin{document}
\begin{center}
Physics - Undergraduate Level Assessment 2 \
Total Marks: 40 \
Answer all questions.
\end{center}

\section*{Multiple Choice (10 marks)}
\begin{enumerate}[label=\arabic*.]
\item In the diamond structure of elemental carbon, the nearest neighbors of each C atom lie at the corners of a
\begin{enumerate}[label=(\alph*)]
    \item square
    \item hexagon
    \item cube
    \item tetrahedron
\end{enumerate}
\item A meter stick with a speed of 0.8 c moves past an observer. In the observer's reference frame, how long does it take the stick to pass the observer ?
\begin{enumerate}[label=(\alph*)]
    \item 1.6 ns
    \item 2.5 ns
    \item 4.2 ns
    \item 6.9 ns
\end{enumerate}
\item For which of the following thermodynamic processes is the increase in the internal energy of an ideal gas equal to the heat added to the gas?
\begin{enumerate}[label=(\alph*)]
    \item Constant temperature
    \item Constant volume
    \item Constant pressure
    \item Adiabatic
\end{enumerate}
\item The speed of light inside of a nonmagnetic dielectric material with a dielectric constant of 4.0 is
\begin{enumerate}[label=(\alph*)]
    \item 1.2 * $10^9$ m/s
    \item 3.0 * $10^8$ m/s
    \item 1.5 * $10^8$ m/s
    \item 1.0 * $10^8$ m/s
\end{enumerate}
\item A rod measures 1.00 m in its rest system. How fast must an observer move parallel to the rod to measure its length to be 0.80 m?
\begin{enumerate}[label=(\alph*)]
    \item 0.50 c
    \item 0.60 c
    \item 0.70 c
    \item 0.80 c
\end{enumerate}
\end{enumerate}

\section*{True / False (10 marks)}
\textit{Answer True or False}
\begin{enumerate}[label=\arabic*.]
\setcounter{enumi}{5}
\item True or False: The speed of light in a nonmagnetic dielectric material with a dielectric constant of 4.0 is 3.0 x $10^8$ m/s.
\item True or False: The detector will detect an average of 10 photons, with an rms deviation of about 3 due to the quantum efficiency of 0.1.
\item True or False: In a reversible thermodynamic process, the temperature of the system must always remain constant.
\item True or False: The primary source of the Sun's energy is a series of thermonuclear reactions involving four hydrogen atoms fusing to form one helium atom.
\item True or False: The particle will travel 450 m in the lab frame before decaying, accounting for relativistic effects.
\end{enumerate}

\section*{Long Form Response (20 marks)}
\textit{Answer each question with a clear and complete explanation.}
\begin{enumerate}[label=\arabic*.]
\setcounter{enumi}{10}
\item Discuss the harmonic series of an organ pipe that is closed at one end and open at the other, specifically analyzing how to determine the frequency of the next higher harmonic given a fundamental frequency of 131 Hz.
\item A police car moving towards a wall at 3.5 m/s emits a siren at 600 Hz. Calculate the frequency of the echo heard by the driver, discussing the Doppler effect involved in your analysis.
\end{enumerate}

\vfill
\noindent\textit{End of Paper}
\end{document}